\section{Avaliação de Ideias e Métricas para Classe Rara (M12)}
\label{sec:m12}

A M12 confrontou \textbf{quatro propostas de extensão} contra a evidência já acumulada pelo projeto,
decidindo o que incorporar e o que rejeitar \emph{com justificativa}. Duas foram implementadas
(ao menos parcialmente) e duas, rejeitadas. O critério foi sempre o mesmo: a complexidade adicional
precisa atacar um gargalo \emph{real} e medido, não apenas conferir verniz arquitetural.

\subsection{Métricas de classe rara: PR-AUC e \emph{lift} (ADR-0015)}
\label{sec:m12-prauc}
Sob forte desbalanceamento (anomalias raras), a \textbf{ROC-AUC engana}: o excesso de
verdadeiros-negativos infla a pontuação. Implementou-se \texttt{src/evaluate.py::score\_curves}
(PR-AUC, ROC-AUC e \emph{baseline}/\emph{lift} sobre o acaso) e adotou-se \textbf{PR-AUC $+$ F1}
como leitura padrão; a ROC-AUC fica apenas como contraste didático. Varrendo a raridade no ground-truth
sintético da PETR4 (Tabela~\ref{tab:m12-prauc}):

\begin{table}[h]
\centering
\caption{PR-AUC \emph{vs.} ROC-AUC ao variar a raridade (ADR-0015, notebook \texttt{14}, PETR4).}
\label{tab:m12-prauc}
\begin{tabular}{rccc}
\toprule
\texttt{n\_injections} & prevalência & \textbf{PR-AUC} & ROC-AUC \\
\midrule
2  & 0,049 & \textbf{0,151} & 0,840 \\
5  & 0,105 & 0,314 & 0,847 \\
10 & 0,188 & 0,537 & 0,921 \\
25 & 0,447 & 0,707 & 0,878 \\
50 & 0,715 & 0,865 & 0,885 \\
\bottomrule
\end{tabular}
\end{table}

\noindent No regime \textbf{raro} (\texttt{n=2}, prevalência 4,9\%) a ROC-AUC $= 0{,}84$ \emph{mascara}
um desempenho fraco que a PR-AUC $= 0{,}15$ expõe (apenas ${\sim}3\times$ o acaso). À medida que a classe
se equilibra, as duas convergem. A crítica é \textbf{confirmada} --- e, de quebra, expôs um defeito do
nosso protocolo sintético.

\subsection{Custo assimétrico FP$\times$FN (ADR-0016, parcial)}
\label{sec:m12-custo}
A escolha de limiar (ADR-0005) é uma \textbf{decisão de risco}, não só estatística: um falso positivo
(alarme falso) e um falso negativo (anomalia perdida) custam diferente. Implementou-se
\texttt{src/evaluate.py::expected\_cost(flags, labels, cost\_fp, cost\_fn)}, que reporta o custo total
sob diferentes razões FN:FP sobre a matriz de confusão. Isso conecta a métrica ao domínio financeiro
sem inventar uma estratégia de \emph{trading}.

\subsection{Ideias rejeitadas com defesa}
\label{sec:m12-rejeitadas}
\begin{itemize}[noitemsep]
  \item \textbf{Atenção / Transformers (ADR-0013):} a premissa (LSTM ``esquece'' em janelas longas)
        \emph{não se aplica} a \texttt{window\_size=30} (ADR-0002); o contexto de longo alcance já entra
        pelo bloco macro (ADR-0012); o dataset ($\sim$2450 janelas/ativo) é pequeno demais para
        Transformers \emph{data-hungry}; e o \texttt{latent\_dim} já se mostrou \textbf{insensível} em
        [8, 32] --- não falta capacidade. Rejeitada.
  \item \textbf{Otimização bayesiana / Optuna (ADR-0014):} a seleção \emph{não} foi manual casual ---
        usou-se walk-forward com \texttt{TimeSeriesSplit} (ADR-0010). A paisagem de \texttt{latent\_dim}
        é \textbf{plana} (diferenças na 4ª casa decimal, abaixo do desvio entre folds $\approx 0{,}003$);
        buscar o ótimo sobre ruído só produz falsa precisão. Rejeitada.
\end{itemize}

\subsection{Trabalho futuro declarado}
\label{sec:m12-futuro}
Dois itens ficam explicitamente como direção futura, e não como entrega:
\begin{itemize}[noitemsep]
  \item \textbf{Backtest de portfólio (ADR-0016):} anomalia $\ne$ sinal de \emph{trade}; um P\&L
        exigiria uma \emph{estratégia} (qual ação tomar?) que o escopo não tem, e um backtest ingênuo
        no período de teste fabricaria facilmente uma história de lucro frágil a \emph{look-ahead} e
        \emph{overfit}. Pré-requisitos claros: definir a ação associada à anomalia, custos de transação,
        protocolo anti-\emph{look-ahead}.
  \item \textbf{Correção do protocolo sintético (consequência da ADR-0015):} com a injeção \emph{default}
        (50 choques) e janelas sobrepostas, cada choque rotula ${\sim}$\texttt{window\_size} janelas e
        a prevalência sobe a ${\sim}$70\% --- a classe \textbf{deixa de ser rara}. Para avaliar no
        regime realista 99/1, reduzir \texttt{n\_injections} e/ou \textbf{agrupar detecções contíguas
        como um único evento} antes de rotular.
\end{itemize}

\noindent\textbf{Limite honesto:} no teste \emph{real} (não supervisionado) não há rótulos, então a
PR-AUC só é computável sobre o sintético; no real, segue-se com a fração marcada $+$ correlação a
eventos (ADR-0008).
